\documentclass[12pt, a4paper, footinclude=true]{scrbook}
\usepackage{geometry}[margin = 1cm]
\usepackage{actuarialsymbol}
\usepackage{actuarialangle}
\usepackage[spanish]{babel}

% Paquetes para dibujo de gráficas 
\usepackage{tikz}
\usepackage{pgfplots}
\usepackage{pgfplotstable}
\usetikzlibrary{arrows.meta, bending, decorations.pathreplacing, shapes.geometric}
\pgfplotsset{compat=1.18}

\usepackage{enumitem}
\usepackage{amsmath}
\usepackage{amssymb}
\usepackage{xcolor}
\usepackage{tcolorbox}
\tcbuselibrary{skins, breakable}
\usetikzlibrary{decorations.pathreplacing}

% Base general
\newtcolorbox{calloutbase}[2][]{
	enhanced,
	breakable,
	sharp corners,
	boxrule=0pt,
	colback=#2!12,
	borderline west={3pt}{0pt}{#2!70!black},
	left=6pt,
	right=6pt,
	top=6pt,
	bottom=6pt,
	#1
}

% Entornos específicos
\newenvironment{callout-azul}{
	\begin{calloutbase}{blue}
	}{
	\end{calloutbase}
}

\newenvironment{advertencia}{
	\begin{calloutbase}{orange}
	}{
	\end{calloutbase}
}

\newenvironment{callout-rojo}{
	\begin{calloutbase}{red}
	}{
	\end{calloutbase}
}

\newenvironment{callout-verde}{
	\begin{calloutbase}{green}
	}{
	\end{calloutbase}
}

\newenvironment{callout-morado}{
	\begin{calloutbase}{violet}
	}{
	\end{calloutbase}
}


% Title Page
\title{Resumen de Matemáticas Actuariales}
\author{Victor Gabriel Vázquez Montes}

\begin{document}

\begin{titlepage}
	
	\huge{Universidad Cristóbal Colón}


\end{titlepage}


\tableofcontents

\chapter{Matemáticas Actuariales I}

\section{Notación probabilística}

``Future Lifetime'' = Vida futura
\begin{center}
\begin{tikzpicture}[scale=1.5]
	
	% Línea principal
	\draw[thick] (0,0) -- (8,0);
	
	% Marcas verticales
	\draw[thick] (0,0.15) -- (0,-0.15);
	\draw[thick] (5,0.15) -- (5,-0.15);
	\draw[thick] (7,0.15) -- (7,-0.15);
	
	% Etiquetas abajo
	\node at (5,-0.5) {$x$};
	\node at (7,-0.5) {$x+tx$};
	
	% Llaves superiores
	\draw[decorate,decoration={brace,amplitude=6pt}]
	(0,0.4) -- (5,0.4)
	node[midway,yshift=0.6cm] {$x$};
	
	\draw[decorate,decoration={brace,amplitude=6pt}]
	(5,0.4) -- (7,0.4)
	node[midway,yshift=0.6cm] {$Fx$};
	
\end{tikzpicture}
\end{center}

Donde: 

\begin{itemize}
\item ($x$) Sera la edad de una persona, por lo que se puede afirmar que $x \geq 0$
\item Se debe tomar en cuenta que la muerte puede ocurrir en cualquier momento mayor a $x$ (No puede ocurrir antes de $x$ porque se supone que la persona ya llegó a esa edad). 
\end{itemize}

\subsection{Glosario}

\begin{itemize}

\item $Fx$: Representa la probabilidad de que ($x$) no sobreviva mas de $x + t$ años. 
\item $Sx$: Es la probabilidad de que una persona sobreviva mas de $x+t$ anos. 
\item $T_x$> Es la cantidad de tiempo hasta que una persona de $x$ años muera. 
\item $T_0$: Vida futura de un recien nacido considerando que el bebe pudo morir antes de llegar a la edad $x$. 
\item 

\subsection{Distribución de la vida futura}

\begin{flushright} $Fx$ para $T_x$ 


\end{flushright}

\begin{tikzpicture}[scale=1.5]
	
	% Línea principal
	\draw[thick] (0,0) -- (8,0);
	
	% Marcas verticales
	\draw[thick] (0,0.15) -- (0,-0.15);
	\draw[thick] (5,0.15) -- (5,-0.15);
	\draw[thick] (7,0.15) -- (7,-0.15);
	
	% Etiquetas abajo
	\node at (0,-0.5) {$0$};
	\node at (5,-0.5) {$t$};
	\node at (7,-0.5) {$\infty$};
	
	% Llaves superiores
	\draw[decorate,decoration={brace,amplitude=6pt}]
	(0,0.4) -- (5,0.4)
	node[midway,yshift=0.6cm] {$Fx = Pr[Tx \leq t]$};
	
\end{tikzpicture}


Siendo así, entonces la función de supervivencia es la probabilidad inversa de que la persona fallezca en dicho periodo, por lo que se puede concluir que: 

$$Sx = 1 - Fx = Pr[Tx < t]$$

\end{itemize}

\section{Notacion Actuarial}

Es diferente la notacion que se emplea en la actuaria a aquella con la que los padres de la probabilidad empezaron a explorar las funciones de supervivencia y mortalidad. 

En actuaria existen notaciones equivalentes a las funciones definidas con anterioridad: 

$$ \px[t]{x} = Pr[T_x > t] = Sx(t)$$

$$ \qx[t]{x} = Pr[T_x < t] = Fx(t)$$

\textbf{Propiedades}

\begin{enumerate}
	\item $\px[t]{x} + \qx[t]{x} = 1$ (O estas muerto o estas vivo). 
	\item $\px[t+u]{x} = \px[t]{x} \cdot \px[u]{x+t}$.
	\subitem Esta propiedad solo es válida para la función de supervivencia y aunque si se puede aplicar a la de fallecimiento, requiere de emplear despejes y ecuaciones algebráicas para obtener el resultado deseado. 
\end{enumerate}

\subsection{Glosario}

\begin{itemize}
	\item $\px[t]{x}$: Es la probabilidad de que una persona sobreviva durante $t$ años teniendo $x$ años. 
	\item $\qx[t]{x}$: Es la probabilidad de que una persona fallezca dentro de los $t$ años siguientes teniendo $x$ años. 
\end{itemize}

\section{Fuerza de Mortalidad}

Es la probabilidad instantanea de fallecer en el proximo instante, condicionada a haber sobrevivido hasta ese momento. 

\begin{itemize}
	\item Se mide en unidades de ``por año'' (Si el tiempo esta en años). Por ejemplo: Si $\mu_{70} = 0.05$ significa que una persona de 70 años cumplidos tiene una probabilidad aproximada del 5\% de morir durante el proximo año si ese valor se mantiene constante durante en ese intervalo muy corto. 
	\item Simplificacion: Mide que tan ``intensa'' es la mortalidad en ese instante exacto de edad. 
\end{itemize}

\begin{tikzpicture}[scale=1.5]
	
	% Línea principal (eje del tiempo)
	\draw[thick] (0,0) -- (9,0);
	
	% Marcas de tiempo
	\draw[thick] (2,0.15) -- (2,-0.15);
	\draw[thick] (6,0.15) -- (6,-0.15);
	\draw[thick] (8,0.15) -- (8,-0.15);
	
	% Etiquetas de tiempo
	\node at (2,-0.6) {$0$};
	\node at (6,-0.6) {$t$};
	\node at (8,-0.6) {$t+\Delta t$};
	
	% Función de supervivencia
	\node at (4,0.4) {$S_x(t)$};
	
	% Llave grande: sobrevive entre 0 y t
	\draw[decorate,decoration={brace,amplitude=8pt}]
	(6,-1.2) -- (2,-1.2)
	node[midway,yshift=-0.8cm] {Sobrevive entre $0$ y $t$};
	
	% Llave pequeña: muerte entre t y t+Δt
	\draw[decorate,decoration={brace,amplitude=6pt}]
	(6,0.5) -- (8,0.5)
	node[midway,yshift=0.7cm] {Muerte};
	
	% Llave grande inferior: muerte ocurre entre t y t+Δt
	\draw[decorate,decoration={brace,amplitude=10pt}]
	(8,-2.4) -- (2,-2.4)
	node[midway,yshift=-0.9cm] {Muerte ocurre entre $t$ y $t+\Delta t$};
	
\end{tikzpicture}

\section{Leyes de probabilidad}

\subsection{De Moivre}

Especifica que la muerte se distribuye de forma uniforme en $[0, w]$, Donde $w$ es la edad limite de una poblacion. 

$$\mu_x = \frac{1}{w-x}, \hspace{1cm} Sx = 1 - \frac{x}{w}, \hspace{1cm} 0 < x \leq w$$

\subsection{Gompertz}

Asumio que $\ln{|\mu_x|}$ en una recta podía ser expresado como: 

$$\ln{|\mu_x|} = \alpha + \beta x$$

Donde si se desea conocer el valor de $\mu_x$ aplicamos la constante de euler> 

\subsection{Makeham}

Es la misma que \textbf{Gompertz} solo que agregando el parametro adicional $A$. 

La interpretación es que $BC^x$ representa la muerte como deterioración, la cual incrementa exponencialmente con $A$ que representa una constante. 

\section{Limites de tiempo}

Nos sirve para sacar el promedio de años que vivirá una persona durante los siguientes $n$ años. 


$n$ - year composite expectation of life: 
$\min(Tx, n)$

$$\eringx{\endowxn} = E[\min(Tx, n)] = \int_{0}^{n} \px[t]{x} dt \hspace{1cm}\text{\color{red} (Para el caso contínuo)}$$ 

$$e_{\endowxn} = \sum_{k=1}^{n} \px[k]{x} \hspace{1cm} \text{\color{red} (Para el caso discreto)}$$

\begin{center}
	\textit{``En vez de decirnos la probabilidad de que una persona muera, nos dice cuando va a morir''}
\end{center}

\textbf{Relación entre la continua y la discreta}

$$\eringx{x} = \frac{1}{2} + e_x$$

\section{Tablas de mortalidad}

Una tabla de mortalidad especifica el número de vidas a una cierta edad $x$. 

El número de vidas a edad $x_0$ es conocido como \textbf{Radix} o ``Población inicial'' 

\subsection{Tablas de mortalidad selecta}

\subsubsection{Curtate Future Lifetime}

Significa \textbf{Tiempo entero de vida}

Es una variable entera/discreta donde $k_x = k$

\textit{``Intervalo de tiempo de muerte''}

$k_x = k$ significa que la muerte de una persona de edad $x$ estará dentro del k-enésimo intervalo $(k-1, k]$

$$k_x = \lfloor T_x \rfloor \hspace{1cm} \text{\color{red} Valor truncado de $x$}$$ 

\section{Seguros}

\subsection{Tipos de seguros}

\subsubsection{Vitalicio / Vida entera (Whole life)}

Lo que hace este tipo de seguro es asumir que se pagará \$1 al momento de la muerte. 

Por lo tanto, lo que se hace es traer el valor presente de \$ 1 pagadero al momento de la muerte. 

Para ello usaremos una $T_x$ continua

\[
	Z = V^{T_x} = e^{-\delta T_x}
\]

Donde: 
\begin{itemize}
	\item $\delta$: Es la fuerza de interés. 
	\item $T_x$: Es el tiempo que debe regresar el  \$ 1 al presente 
\end{itemize}

Por lo tanto, se puede decir que: 

\begin{itemize}
	\item $E[T_x]$: Es el valor esperado de cuando ocurrirá la muerte. 
	\item $E[Z]$: Es el valor esperado de cuando esperamos pagar la suma asegurada.
\end{itemize}

\[
	E[Z] = E[e^{-\delta T_x}]
\]


\begin{center}
\tikzset{every picture/.style={line width=0.75pt}} %set default line width to 0.75pt        
\begin{tikzpicture}[x=0.75pt,y=0.75pt,yscale=-1,xscale=1]
	%uncomment if require: \path (0,300); %set diagram left start at 0, and has height of 300
	
	%Straight Lines [id:da7376839101373023] 
	\draw    (109,161) -- (527.33,161) (273,157) -- (273,165)(437,157) -- (437,165) ;
	\draw [shift={(530.33,161)}, rotate = 180] [fill={rgb, 255:red, 0; green, 0; blue, 0 }  ][line width=0.08]  [draw opacity=0] (8.93,-4.29) -- (0,0) -- (8.93,4.29) -- cycle    ;
	\draw [shift={(109,161)}, rotate = 180] [color={rgb, 255:red, 0; green, 0; blue, 0 }  ][line width=0.75]    (0,5.59) -- (0,-5.59)   ;
	%Shape: Brace [id:dp502356150507335] 
	\draw   (108,199) .. controls (108,203.67) and (110.33,206) .. (115,206) -- (180.17,206) .. controls (186.84,206) and (190.17,208.33) .. (190.17,213) .. controls (190.17,208.33) and (193.5,206) .. (200.17,206)(197.17,206) -- (265.33,206) .. controls (270,206) and (272.33,203.67) .. (272.33,199) ;
	%Shape: Rectangle [id:dp2116147331087268] 
	\draw  [color={rgb, 255:red, 0; green, 0; blue, 0 }  ,draw opacity=0 ][fill={rgb, 255:red, 208; green, 2; blue, 27 }  ,fill opacity=0.3 ] (273,154) -- (436.33,154) -- (436.33,170) -- (273,170) -- cycle ;
	%Shape: Free Drawing [id:dp8495024415521517] 
	\draw  [color={rgb, 255:red, 184; green, 233; blue, 134 }  ,draw opacity=1 ][line width=3] [line join = round][line cap = round] (176.33,89.67) .. controls (176.33,84.61) and (167.54,88.46) .. (163.67,92.33) .. controls (157.56,98.44) and (156.92,120.96) .. (166.33,126.33) .. controls (174.31,130.89) and (176.02,130.49) .. (186.33,129.67) .. controls (213.77,127.47) and (207.22,70.89) .. (177.67,85.67) ;
	%Shape: Free Drawing [id:dp7818874018128433] 
	\draw  [color={rgb, 255:red, 184; green, 233; blue, 134 }  ,draw opacity=1 ][line width=3] [line join = round][line cap = round] (173.67,98.33) .. controls (173.67,102.22) and (174.33,112.94) .. (174.33,109) ;
	%Shape: Free Drawing [id:dp37608410612369947] 
	\draw  [color={rgb, 255:red, 184; green, 233; blue, 134 }  ,draw opacity=1 ][line width=3] [line join = round][line cap = round] (182.33,97) .. controls (182.33,100.34) and (182.34,103.72) .. (183,107) ;
	%Shape: Free Drawing [id:dp6273715752908713] 
	\draw  [color={rgb, 255:red, 184; green, 233; blue, 134 }  ,draw opacity=1 ][line width=3] [line join = round][line cap = round] (168.33,111) .. controls (168.33,125.06) and (188.33,123.34) .. (188.33,109.67) ;
	%Shape: Free Drawing [id:dp715308403908088] 
	\draw  [color={rgb, 255:red, 208; green, 2; blue, 27 }  ,draw opacity=0.3 ][line width=3] [line join = round][line cap = round] (347,191) .. controls (335.63,191) and (330.05,209.48) .. (331,219) .. controls (332.79,236.9) and (356.59,235.89) .. (368.33,231) .. controls (379.99,226.14) and (378.74,198.27) .. (370.33,191.67) .. controls (364.67,187.21) and (350,187.87) .. (344.33,189) ;
	%Shape: Free Drawing [id:dp285519545736307] 
	\draw  [color={rgb, 255:red, 208; green, 2; blue, 27 }  ,draw opacity=0.3 ][line width=3] [line join = round][line cap = round] (352.33,199.67) .. controls (349.67,203.82) and (347.16,208.18) .. (343.67,211.67) ;
	%Shape: Free Drawing [id:dp23706621271126382] 
	\draw  [color={rgb, 255:red, 208; green, 2; blue, 27 }  ,draw opacity=0.3 ][line width=3] [line join = round][line cap = round] (351.67,210.33) .. controls (350.71,208.42) and (343.4,201.26) .. (342.33,202.33) ;
	%Shape: Free Drawing [id:dp2615874294505175] 
	\draw  [color={rgb, 255:red, 208; green, 2; blue, 27 }  ,draw opacity=0.3 ][line width=3] [line join = round][line cap = round] (365.67,197) .. controls (363.89,200.56) and (364.31,207.67) .. (360.33,207.67) ;
	%Shape: Free Drawing [id:dp7482425513540865] 
	\draw  [color={rgb, 255:red, 208; green, 2; blue, 27 }  ,draw opacity=0.3 ][line width=3] [line join = round][line cap = round] (367,207) .. controls (363.48,205.24) and (362.79,197.67) .. (358.33,197.67) ;
	%Shape: Free Drawing [id:dp6228508298661153] 
	\draw  [color={rgb, 255:red, 208; green, 2; blue, 27 }  ,draw opacity=0.3 ][line width=3] [line join = round][line cap = round] (365,220.33) .. controls (362.37,217.7) and (358.01,216.42) .. (354.33,217) .. controls (350.98,217.53) and (350.15,219.85) .. (348.33,221.67) .. controls (347.94,222.06) and (347,224.74) .. (347,223) ;
	
	% Text Node
	\draw (104,135.4) node [anchor=north west][inner sep=0.75pt]    {$0$};
	% Text Node
	\draw (104,170.4) node [anchor=north west][inner sep=0.75pt]    {$x$};
	% Text Node
	\draw (255,170.4) node [anchor=north west][inner sep=0.75pt]    {$x+s$};
	% Text Node
	\draw (269,136.4) node [anchor=north west][inner sep=0.75pt]    {$s$};
	% Text Node
	\draw (421,136.4) node [anchor=north west][inner sep=0.75pt]    {$s+\triangle s$};
	% Text Node
	\draw (171,224.4) node [anchor=north west][inner sep=0.75pt]    {$\px[s]{x}$};
	% Text Node
	\draw (397,169.4) node [anchor=north west][inner sep=0.75pt]    {$x+s+\triangle \delta s\ $};
	% Text Node
	\draw (336,129.4) node [anchor=north west][inner sep=0.75pt]  [color={rgb, 255:red, 208; green, 2; blue, 27 }  ,opacity=0.5 ]  {$\mu _{x+s}$};
	% Text Node
	\draw (279,108) node [anchor=north west][inner sep=0.75pt]  [color={rgb, 255:red, 208; green, 2; blue, 27 }  ,opacity=0.5 ] [align=left] {Fuerza de mortalidad};
\end{tikzpicture}
\end{center}

Ilustrado el concepto, entonces tenemos que explicar que: 

\begin{itemize}
	\item $x$: La edad a la que una persona contrata el seguro. 
	\item $x+s$: Llega con vida
	\item $\mu_{x+s}$: Muere
\end{itemize}

\[
	\Ax*{x} = E[Z] = E[e^{-\delta T_x}] = \int_{0}^{\infty} e^{-\delta t} \cdot \px[t]{x} \cdot \mu_{x+t} dt
\]

$\Ax*{x}$ es entonces el valor presente actuarial (VPA) de un beneficio (Suma asegurada) para una persona en edad $x$ pagadero al momento de la muerte. 

Donde la suma asegurada es igual a $S$ por lo tanto $S \rightarrow E[S \cdot Z] = S \cdot \Ax*{x}$

{\color{red} Si tenemos fuerza de mortalidad constante, entonces:}

\[
	E[Z] = \int_{0}^{\infty} e^{-\delta t} \cdot \px[t]{x} \cdot \mu_{x+t} dt \\ 
\]

\[
	\text{Se debe recordar que...} \hspace{0.5cm} \underbrace{\px[t]{x} = e^{\int_{0}^{y} \mu x dx}}_{\color{red}Continua}
\]		

\[
	\text{Por lo que, dado que $\mu_x$ es constante, entonces tenemos...} \hspace{0.5cm} \px[t]{x} = e^{-\mu t}
\]

De donde obtenemos que: 

\[
	\Ax*{x} = \int_{0}^{\infty} e^{-\delta t} \cdot e^{-\mu t} \cdot \mu dt
\]

\[
	\mu \int_{0}^{\infty} e^{-\delta t - \mu t} dt = \mu \cdot e^{(\delta + \mu) |_{0}^{\infty}} \cdot \frac{1}{-(\mu + \delta)}
\]

\[
	-\frac{\mu}{\mu + \delta} e^{- \infty} + \frac{\mu}{\mu + \delta} e^{-0} = \boxed{\color{red} \frac{\mu}{\mu + \delta}}
\]

\subsubsection{Seguro Temporal (Term Insurance)}

\subsubsection{Seguro dotal puro (Pure Endowment)}

Es el tipo de seguro que se paga si y solo si el asegurado llega con vida a los $n$ años. 

$$
	Z = 
	\begin{cases}
		0 & \text{Con probabilidad $\qx[n]{x}$} \\
		V^n = e^{-\delta n} & \text{Con probabilidad $\px[n]{x}$} 
	\end{cases}
$$

Por lo tanto su notación es comprendida de la siguiente manera: 

$$ \Ax*{\pureendowxn} = \Ex[n]{x}$$

$$E[Z] = \Ax*{\pureendowxn} = v^n \cdot \px[n]{x} = e^{-\delta n} \cdot \px[n]{x}$$

Donde, si $\px[n]{x}$ tiene una mortalidad constante (Fuerza de mortalidad constante):

$$\Ax*{\pureendowxn} = e^{-\delta n} \cdot e^{-\mu n} = e^{-n(\delta + \mu)}$$

\subsubsection{Seguro dotal mixto (Endowment)}

\[
	Z = 
	\begin{cases}
		v^{T_x} = e^{-\delta T_x} & T_x < n \rightarrow \text{Si muere} \\
		
		v^n = e^{-\delta n} & T_x \geq n \rightarrow \text{Si vive}
	\end{cases}
\]

\[
	Z = V^{\min(T_x, n)} = e^{-\delta \min(T_x, n)}
\]

Por lo que: 

\[
	\Ax*{\endowxn} = \Ax*{\termxn} + \Ax{\pureendowxn}	= \left( \int_{0}^{n} e^{-\delta t} \cdot \px[t]{x} \cdot \mu_{x+t} \delta t\right) + \left( V_{t}^{n} \px{x} \right)
\]

El seguro dotal mixto resulta en una combinación entre un seguro temporal y un seguro dotal puro, debido a que paga tanto si la persona sobrevive como si la persona muere. 

\subsection{Glosario y Notación}

$\Ax{x}$: Valor presente actuarial (VPA) o en inglés (NPV) de una unidad monetaria pagadera al final del año de fallecimiento de una persona de edad $x$ (Whole life insurance).\\

$\Ax{\termxn}$: Denota el valor presente actuarial de un seguro temporal (n-year term insurance). 

$\Ax{x:\angln}$: Denota el valor presente actuarial de un seguro dotal mixto (n-year endowment insurance)

$\Ax{\pureendowxn}$: Denota el valor presente actuarial (VPA) de un seguro dotal puro (n-year pure endowment)

$\Ex[n]{x}$: Factor de descuento actuarial (The pure endowment present value factor)

$\Ax{x}[(m)]$: Denota el valor presente actuarial para un seguro en el que el pago por la muerte se efectúa al final del m-ésimo periodo del año en el que ocurre la muerte. 

\begin{itemize}
	\item $(m)$: Es la frecuencia de pagos (el seguro se paga en fracciones del año).
	\item Es un seguro de vida entera (cubre toda la vida de la persona).
	\item Donde $(m) > 0$ por lo que: 
	\subitem $m = 2$ Semestral cada 6 meses = Dos veces al año "semi-annual". 
	\subitem $m = 4$ Trimestral cada 3 meses = 4 veces al año "quarterly".
	\subitem $m = 12$ Mensual cada mes, 12 veces al año "monthly". 
	\subitem $m = 6$ bimestral cada 2 meses = 6 veces al año.
	\subitem La lista puede continuar dependiendo cuantas veces y que tan pagadero se quiera hacer el seguro.  
\end{itemize}

Como comentario, 

\section{Anualidades contingentes}

Anualidad contingente: Serie de pagos al final de un periodo condicionada a la vida de una persona

\begin{itemize}
	\item Puede ser empleado para seguros o pensiones que empiezan a partir de cierto periodo y terminan hasta que el pensionado o asegurado muere. 
	\item Incluso podría llegar a ser aplicado para alguna segubeca o seguros similares. 
\end{itemize}

\section{Edades fraccionarias de mortalidad}

\textbf{UDD}: Uniform Distribution of death. 

\textbf{Se asume una fuerza constante de mortalidad}

\vspace{1cm}

\tikzset{every picture/.style={line width=0.75pt}} %set default line width to 0.75pt 

\begin{tikzpicture}[x=0.75pt,y=0.75pt,yscale=-1,xscale=1]
	%uncomment if require: \path (0,300); %set diagram left start at 0, and has height of 300
	
	%Straight Lines [id:da23710618032520359] 
	\draw    (104,117) -- (480,117) (284,110.5) -- (284,123.5)(464,110.5) -- (464,123.5) ;
	\draw [shift={(483,117)}, rotate = 180] [fill={rgb, 255:red, 0; green, 0; blue, 0 }  ][line width=0.08]  [draw opacity=0] (8.93,-4.29) -- (0,0) -- (8.93,4.29) -- cycle    ;
	\draw [shift={(101,117)}, rotate = 0] [fill={rgb, 255:red, 0; green, 0; blue, 0 }  ][line width=0.08]  [draw opacity=0] (8.93,-4.29) -- (0,0) -- (8.93,4.29) -- cycle    ;
	%Straight Lines [id:da3667942171299532] 
	\draw    (103,50) -- (321,49) ;
	%Straight Lines [id:da07887529438472474] 
	\draw    (321,118) -- (321,49) ;
	%Straight Lines [id:da42361193166684863] 
	\draw  [dash pattern={on 0.84pt off 2.51pt}]  (284,119) -- (284,50) ;
	%Shape: Rectangle [id:dp2856002199031109] 
	\draw  [color={rgb, 255:red, 255; green, 0; blue, 0 }  ,draw opacity=0.2 ][fill={rgb, 255:red, 255; green, 0; blue, 0 }  ,fill opacity=0.2 ] (284.01,117.79) -- (284.39,49) -- (321.38,49.21) -- (321,118) -- cycle ;
	%Straight Lines [id:da23734398010970692] 
	\draw [color={rgb, 255:red, 255; green, 0; blue, 0 }  ,draw opacity=0.2 ]   (287,130) -- (317,130) ;
	\draw [shift={(320,130)}, rotate = 180] [fill={rgb, 255:red, 255; green, 0; blue, 0 }  ,fill opacity=0.2 ][line width=0.08]  [draw opacity=0] (8.93,-4.29) -- (0,0) -- (8.93,4.29) -- cycle    ;
	\draw [shift={(284,130)}, rotate = 0] [fill={rgb, 255:red, 255; green, 0; blue, 0 }  ,fill opacity=0.2 ][line width=0.08]  [draw opacity=0] (8.93,-4.29) -- (0,0) -- (8.93,4.29) -- cycle    ;
	
	% Text Node
	\draw (276,135.4) node [anchor=north west][inner sep=0.75pt]    {$kx$};
	% Text Node
	\draw (441,135.4) node [anchor=north west][inner sep=0.75pt]    {$kx+1$};
	% Text Node
	\draw (104,136.4) node [anchor=north west][inner sep=0.75pt]    {$0$};
	% Text Node
	\draw (50,138.4) node [anchor=north west][inner sep=0.75pt]    {$t$};
	% Text Node
	\draw (83,181) node [anchor=north west][inner sep=0.75pt]   [align=left] {Edad $\displaystyle x$};
	
	
\end{tikzpicture}

\chapter{Matemáticas Actuariales II}

\section{Reservas}

La reserva matemática de una póliza en vigor, al momento de valuación se determinará con base en la diferencia entre el VPA de obligaciones y el VPA de las primas

Se asumirá siempre: 

\begin{itemize}
	\item $i = 0.06$ Tabla de mortalidad de la SOA
\end{itemize}

La reserva a tiempo $t$ será estimada como el VP de la pérdida futura a tiempo $t$. 

$$t_V = E[L_t]$$

\begin{center}
\tikzset{every picture/.style={line width=0.75pt}} %set default line width to 0.75pt        

\begin{tikzpicture}[x=0.75pt,y=0.75pt,yscale=-1,xscale=1]
	%uncomment if require: \path (0,300); %set diagram left start at 0, and has height of 300
	
	%Curve Lines [id:da31604371037346146] 
	\draw    (282.33,72) .. controls (215.35,68.06) and (195.91,95.17) .. (200.12,136.12) ;
	\draw [shift={(200.33,138)}, rotate = 263.21] [color={rgb, 255:red, 0; green, 0; blue, 0 }  ][line width=0.75]    (10.93,-3.29) .. controls (6.95,-1.4) and (3.31,-0.3) .. (0,0) .. controls (3.31,0.3) and (6.95,1.4) .. (10.93,3.29)   ;
	%Shape: Boxed Bezier Curve [id:dp600694142508993] 
	\draw    (364.42,232.27) .. controls (425.51,239.13) and (423.4,202.22) .. (426.08,164.63) ;
	\draw [shift={(426.2,162.91)}, rotate = 94.31] [color={rgb, 255:red, 0; green, 0; blue, 0 }  ][line width=0.75]    (10.93,-3.29) .. controls (6.95,-1.4) and (3.31,-0.3) .. (0,0) .. controls (3.31,0.3) and (6.95,1.4) .. (10.93,3.29)   ;
	%Shape: Boxed Bezier Curve [id:dp8151678362789523] 
	\draw    (199.34,170) .. controls (191.19,231.07) and (233.89,229.08) .. (277.35,231.87) ;
	\draw [shift={(279.33,232)}, rotate = 183.89] [color={rgb, 255:red, 0; green, 0; blue, 0 }  ][line width=0.75]    (10.93,-3.29) .. controls (6.95,-1.4) and (3.31,-0.3) .. (0,0) .. controls (3.31,0.3) and (6.95,1.4) .. (10.93,3.29)   ;
	%Shape: Boxed Bezier Curve [id:dp1741955785108219] 
	\draw    (424.77,134.69) .. controls (431.21,73.55) and (394.31,75.91) .. (356.7,73.49) ;
	\draw [shift={(354.99,73.38)}, rotate = 3.92] [color={rgb, 255:red, 0; green, 0; blue, 0 }  ][line width=0.75]    (10.93,-3.29) .. controls (6.95,-1.4) and (3.31,-0.3) .. (0,0) .. controls (3.31,0.3) and (6.95,1.4) .. (10.93,3.29)   ;
	
	% Text Node
	\draw (297,141) node [anchor=north west][inner sep=0.75pt]   [align=left] {{\fontfamily{ptm}\selectfont Seguro}};
	% Text Node
	\draw (300,65) node [anchor=north west][inner sep=0.75pt]   [align=left] {{\fontfamily{ptm}\selectfont Prima}};
	% Text Node
	\draw (293,221) node [anchor=north west][inner sep=0.75pt]   [align=left] {{\fontfamily{ptm}\selectfont Beneficio}};
	% Text Node
	\draw (160,142) node [anchor=north west][inner sep=0.75pt]   [align=left] {{\fontfamily{ptm}\selectfont Aseguradora}};
	% Text Node
	\draw (393,141) node [anchor=north west][inner sep=0.75pt]   [align=left] {{\fontfamily{ptm}\selectfont Asegurado}};
	
\end{tikzpicture}
	
\end{center}

La ilustración anterior sirve para recordar de forma práctica lo que es un seguro y su naturaleza cíclica, en donde un asegurado paga una prima, que ingresa a una aseguradora para posteriormente pagar un beneficio a ese asegurado.\\ 

Es además importante recordar los tipos de primas que existen: 

\begin{itemize}
	\item Prima única: Aquella en la que se efectúa un solo pago al principio. 
	\item Prima nivelada: Aquella en la que se efectúan una serie de pagos a lo largo de un periodo. 
\end{itemize}

\textbf{El principio de equivalencia:} Sostiene que el valor actuarial esperado de las primas (aportaciones), calculados ambos en el mismo instante del tiempo y bajo las mismas hipótesis financieras y biométricas es igual al valor actuarial esperado de las primas, en otras palabras: 

$$ \text{VAE(Beneficios)} = \text{VAE(Primas)}$$

Esto es de suma importancia, debido a que implica que el contrato de seguros debe ser financieramente justo (equilibrado) entre ambas partes, lo cuál se cumple si: 

\begin{itemize}
	\item La aseguradora no obtiene ganancia ni pérdida neta. 
	\item El asegurado no subsidia ni es subsididado.
\end{itemize}

Este principio es meramente técnico al implicar neutralidad financiera, puesto que las aseguradoras (como toda entidad lucrativa) tienen el mantra de tener rentabilidad. 

Para que el principio sea válido, se consideran los siguientes elementos: 

\begin{itemize}
	\item Interés técnico $i$ 
	\item Mortalidad/Supervivencia (tablas biométricas)
	\item Probabilidades de ocurrencia del siniestro
	\item Estructura temporal de pagos
	\item Valor del dinero en el tiempo
\end{itemize}

$l$: Es como se denota la pérdida futura en la que va a incurrir la compañía aseguradora. 

\begin{enumerate}
	\item Prima neta ``Net''
	
	$$ L_0^n = \text{VP de Beneficios} - \text{VP de las primas} \hspace{.5cm} \text{(Sin gastos)}$$
	
	\item Prima bruta ``Gross''
	
	\[
		L_0^n = \text{VP beneficios} + \text{VP de los gastos} - \text{VP de las primas} \hspace{.5cm} \text{(Con gastos)}
	\]
\end{enumerate} 

\textbf{Primas}

Usando el principio de equivalencia


$$E[L_0^n] = 0 \hspace{.5cm} \text{``El valor esperado de la pérdida futura será 0''}$$

\textbf{Ejemplo}

\vspace{.5cm}

Un asegurado de (60) compra un seguro vitalicio de \$ 100,000 pagadero al momento de su muerte; Se nos encarga estimar la prima bajo el principio de equivalencia, tomando en cuenta que las primas se pagarán hasta que cumpla 80 años, dado que:

\[
	\mu = 0\decimalpoint.05, \hspace{0.5cm} \delta = 0\decimalpoint.006
\] 

\[
	L_0^n = \underbrace{\text{VP Obligaciones}}_{\text{Seguro}} - \text{VP de las primas}
\]

Lo cuál quedaría de la siguiente manera en notación actuarial:

\[
	100,000 \Ax*{60} - P \cdot \ddot{a}_{\endow{60}{20}}
\]

Aplicando el principio de equivalencia, tenemos que la expresión anterior equivale a: 

\[
	0 = 100,000 \Ax*{60} - P \cdot \ddot{a}_{\endow{60}{20}}
\]

Aplicando despejes algebráicos simples llegamos a la conclusión de que: 

\[
	P = \frac{100,000 \Ax*{60}}{\ddot{a}_{\endow{60}{20}}}
\]

Acto seguido calculamos los dos elementos que necesitamos para obtener el valor de nuestra prima: 

\[
	\Ax*{60} = \frac{\mu}{\mu + \delta} = \frac{0.05}{0.056} = 0 \decimalpoint. 8928
\]

Para la parte de la anualidad contingente es necesario recordar que la tasa de descuento $d = 1 - e^{-\delta}$, por lo que: 

\[
	% Primera parte de la expresion
	d =  1 - e ^ {-0\decimalpoint.006} = \decimalpoint.00598203
\]	

\[
	\ddot{a}_{\endow{60}{20}} = \frac{1 - \Ax*{x}{n}}{d} = \frac{1 - [1 - e^{-20(0\decimalpoint.05)(0\decimalpoint.006)}]}{\decimalpoint.00598203} = 54\decimalpoint.54332302 \approx 54\decimalpoint.54 \%
\]

\[
	= \frac{100,000[.8928]}{54.54} = \$ 1638.863965 \approx \boxed{\$ 1638.83}
\]

\textbf{Gross Premium}

Gastos

\begin{itemize}
	\item Gasto inicial: Comisiones del agente + ``Underwrittings''.
	\item Gastos administrativos
	\item Gastos al momento del beneficio
\end{itemize}

\textbf{Ejemplo}

Se tiene un seguro de (30) de vida entera con una suma asegurada de \$ 500,000 pagadero al final del mes, donde las primas serán pagadas al principio de cada año durante el tiempo de cobertura. 

Algunos de los gastos son: 

\begin{enumerate}
	\item 50\% de la primera prima para el agente  \item \$ 2,500 por emisión 
	\item 5\% de cada prima por admin. 
	\item 4\% cuando se cobre la S.A
\end{enumerate}

\section{Vidas Múltiples}

Implica hacer todos los cálculos que hicimos en este semestre y en el anterior, pero haciendolos para muchas personas. 

Se usa para seguros un poco más complejos, como los seguros colectivos, para grupos relacionados entre si, o seguros familiares donde la vida depende de algún familiar. 

Dicho esto ocurren dos casos: 

\begin{itemize}
	\item Vida conjunta (``Primero que muere'')
	\item Último sobreviviente (``Último que muere'')
\end{itemize}

\subsection{Vida conjunta}

Teniendo a dos personas que serán denotadas como $(x)$ y $(y)$, el estatus que se les otorga de ``vida conjunta'' estará vigente siempre y cuando $(x)$ y $(y)$ estén con vida (La misma restricción aplica para un caso de tres o más personas), donde dicho estatus termina o vence en el momento en el que la primera persona muere sin importar quien sea.\\ 

%$(x)$ y $(y)$ son dos personas. \\

%El estatus de vida conjunta estará vigente siempre y cuando $(x)$ y $(y)$ estén con vida (Esta restricción aplica tanto para un caso de dos o más personas). 

%El estatus de vida conjunta se termina/vence cuando la primera persona muere sin importar quien sea. 

$T_{xy}$: Es el tiempo futuro del estatus vida conjunta osea en el tiempo en el que la primer persona del grupo muere. 

$$T_{xy} = \min(T_x, T_y) \leq T_x \hspace{1cm} \& \hspace{1cm} T_{xy} = \min(T_x, T_y) \leq T_y$$

$$\text{Para el caso de cuatro variables la notación sería...} \hspace{0.5cm} T_{xyzw} = \min{T_x, T_y, T_z, T_w}$$\\

\textbf{Probabilidad de fracaso/éxito}\\

La probabilidad de que el estatus sobreviva viene dada por: 

$$P[(x) \text{ y } (y) \text{ estén con vida } t \text{ años } ] = P[T_{xy} > t] = \px[t]{xy}$$ \\

Mientras que la probabilidad de que el estatus falle (Se debe recordar que es la probabilidad de que una persona de las tantas del grupo muera) viene dada por: 

$$P[(x) \text{ y } (y) \text{ No estén con vida en } t \text{ años } ] = Pr[T_{xy} \leq t] = \qx[t]{xy}$$ \\

Y dado que asumimos vidas independientes:

$$\px[t]{xy} = Pr[T_{xy} > T_x \hspace{.3cm} \& \hspace{.3cm} T_{xy} > T_y] = \px[t]{x} \cdot \px[t]{y}$$

$$\qx[t]{xy} = 1 - \px[t]{xy} = 1 - \px[t]{x} \cdot \px[t]{y} = 1 - (1 - \qx[t]{x})(1-\qx[t]{y})$$

\begin{advertencia}
	No debes confundirte en la notación, recuerda que al igual que en la notación probabilística: $$\qx[t]{xy} \neq \qx[t]{x} \cdot \qx[t]{y}$$
\end{advertencia}



\textbf{Ejemplo:}\\

Teniendo la siguiente tabla de mortalidad, donde $(x)$ es el marido de edad 65 y $(y)$ la mujer de edad 60. se nos pide calcular la probabilidad de que el estatus termine en los siguientes 3 años: \\

\begin{center}
\begin{tabular}{c|c|c|c}
	
	$x$ & $\lx{x}$ & $y$ & $\lx{y}$ \\
	\hline
	65 & 4.302 & 60 & 4.726 \\
	66 & 4.189 & 61 & 4.110 \\
	67 & 4.205 & 62 & 4.127 \\
	68 & 4.219 & 63 & 4.143 \\
	69 & 4.234 & 64 & 4.159 \\
	
\end{tabular}
\end{center}

\textbf{Resolución}\\



\tikzset{every picture/.style={line width=0.75pt}} %set default line width to 0.75pt        
\begin{center}
	\tikzset{every picture/.style={line width=0.75pt}} %set default line width to 0.75pt        
	\begin{tikzpicture}[x=0.75pt,y=0.75pt,yscale=-1,xscale=1]
		%uncomment if require: \path (0,300); %set diagram left start at 0, and has height of 300
		
		%Straight Lines [id:da8697542673867815] 
		\draw    (100,156) -- (523.33,153) (245.97,150.97) -- (246.02,158.97)(391.96,149.93) -- (392.02,157.93) ;
		\draw [shift={(523.33,153)}, rotate = 179.59] [color={rgb, 255:red, 0; green, 0; blue, 0 }  ][line width=0.75]    (0,5.59) -- (0,-5.59)   ;
		\draw [shift={(100,156)}, rotate = 179.59] [color={rgb, 255:red, 0; green, 0; blue, 0 }  ][line width=0.75]    (0,5.59) -- (0,-5.59)   ;
		%Straight Lines [id:da13644970534823986] 
		\draw [color={rgb, 255:red, 0; green, 0; blue, 0 }  ,draw opacity=1 ][fill={rgb, 255:red, 208; green, 2; blue, 27 }  ,fill opacity=1 ]   (100,136) -- (523.33,133) (245.97,130.97) -- (246.02,138.97)(391.96,129.93) -- (392.02,137.93) ;
		%Straight Lines [id:da9763051715174768] 
		\draw [color={rgb, 255:red, 0; green, 0; blue, 0 }  ,draw opacity=1 ][fill={rgb, 255:red, 208; green, 2; blue, 27 }  ,fill opacity=1 ]   (100,101) -- (523.33,98) (245.97,95.97) -- (246.02,103.97)(391.96,94.93) -- (392.02,102.93) ;
		
		% Text Node
		\draw (94,170.4) node [anchor=north west][inner sep=0.75pt]    {$0$};
		% Text Node
		\draw (519,163.4) node [anchor=north west][inner sep=0.75pt]    {$3$};
		% Text Node
		\draw (386,163.4) node [anchor=north west][inner sep=0.75pt]    {$2$};
		% Text Node
		\draw (240,163.4) node [anchor=north west][inner sep=0.75pt]    {$1$};
		% Text Node
		\draw (92,112.4) node [anchor=north west][inner sep=0.75pt]  [color={rgb, 255:red, 208; green, 2; blue, 27 }  ,opacity=1 ]  {$65$};
		% Text Node
		\draw (237,111.4) node [anchor=north west][inner sep=0.75pt]  [color={rgb, 255:red, 208; green, 2; blue, 27 }  ,opacity=1 ]  {$66$};
		% Text Node
		\draw (382,111.4) node [anchor=north west][inner sep=0.75pt]  [color={rgb, 255:red, 208; green, 2; blue, 27 }  ,opacity=1 ]  {$67$};
		% Text Node
		\draw (513,112.4) node [anchor=north west][inner sep=0.75pt]  [color={rgb, 255:red, 208; green, 2; blue, 27 }  ,opacity=1 ]  {$68$};
		% Text Node
		\draw (56,111.4) node [anchor=north west][inner sep=0.75pt]  [color={rgb, 255:red, 208; green, 2; blue, 27 }  ,opacity=1 ]  {$x$};
		% Text Node
		\draw (94,77.4) node [anchor=north west][inner sep=0.75pt]  [color={rgb, 255:red, 74; green, 144; blue, 226 }  ,opacity=1 ]  {$60$};
		% Text Node
		\draw (238,76.4) node [anchor=north west][inner sep=0.75pt]  [color={rgb, 255:red, 74; green, 144; blue, 226 }  ,opacity=1 ]  {$61$};
		% Text Node
		\draw (384,76.4) node [anchor=north west][inner sep=0.75pt]  [color={rgb, 255:red, 74; green, 144; blue, 226 }  ,opacity=1 ]  {$62$};
		% Text Node
		\draw (515,77.4) node [anchor=north west][inner sep=0.75pt]  [color={rgb, 255:red, 74; green, 144; blue, 226 }  ,opacity=1 ]  {$63$};
		% Text Node
		\draw (58,76.4) node [anchor=north west][inner sep=0.75pt]  [color={rgb, 255:red, 74; green, 144; blue, 226 }  ,opacity=1 ]  {$y$};
		
		
	\end{tikzpicture}
\end{center}

$$\qx[3]{65:60} = 1 - \px[3]{65:60} = 1 - \px[3]{65} \cdot \px[3]{60} = 1 - \left( \frac{\lx{68}}{\lx{65}}\right) \left(\frac{\lx{63}}{\lx{60}} \right)$$

$$= 1 - \left( \frac{41.351}{43.302}\right) \left(\frac{46.580}{47.260} \right) = 0.06041 = \boxed{6.041} \%$$

\subsection{Último sobreviviente}

Estatus que falla cuando {\color{red} TODAS} las personas del grupo mueren. \\


$$T_{\joint{xy}} = \max(T_x, T_y)\hspace{1cm} \text{Tiempo futuro de vida del grupo}$$ 

$$\px[t]{\joint{xy}} = Pr[ \text{ al menos } (x) \text{ o } (y) \text{ esten con vida } t \text{ años } ] =Pr[T_{\joint{xy}} > t]$$

$$\qx[t]{\joint{xy}} = Pr[ \text{ al menos } (x) \text{ y } (y) \text{ mueran en } t \text{ años } ] =Pr[T_{\joint{xy}} \leq t]
$$\\



\textbf{Relación}

\begin{enumerate}
	\item Si $(x)$ muere primero $\rightarrow T_x = T_{xy} \therefore T_y = T_{\joint{xy}}$
	\item Si $(y)$ muere primero $\rightarrow T_y = T_{xy} \therefore T_x = T_{\joint{xy}}$
\end{enumerate}

\vspace{1cm}

\tikzset{every picture/.style={line width=0.75pt}} %set default line width to 0.75pt        

\begin{tikzpicture}[x=0.75pt,y=0.75pt,yscale=-1,xscale=1]
	%uncomment if require: \path (0,300); %set diagram left start at 0, and has height of 300
	
	%Straight Lines [id:da4363365428126367] 
	\draw    (46,125) -- (486.33,122.67) (195.98,120.21) -- (196.02,128.21)(345.97,119.41) -- (346.02,127.41) ;
	\draw [shift={(486.33,122.67)}, rotate = 179.7] [color={rgb, 255:red, 0; green, 0; blue, 0 }  ][line width=0.75]    (0,5.59) -- (0,-5.59)   ;
	\draw [shift={(46,125)}, rotate = 179.7] [color={rgb, 255:red, 0; green, 0; blue, 0 }  ][line width=0.75]    (0,5.59) -- (0,-5.59)   ;
	%Straight Lines [id:da5570077513947393] 
	\draw  [dash pattern={on 0.84pt off 2.51pt}]  (345.33,160.67) -- (345.33,201.67) ;
	%Straight Lines [id:da9209038411868883] 
	\draw  [dash pattern={on 0.84pt off 2.51pt}]  (194.33,160.67) -- (194.33,175.67) ;
	
	% Text Node
	\draw (160,96.4) node [anchor=north west][inner sep=0.75pt]    {$T_{x} \ =\ T_{xy}$};
	% Text Node
	\draw (307,95.4) node [anchor=north west][inner sep=0.75pt]    {$T_{y} \ =\ T_{\overline{xy}}$};
	% Text Node
	\draw (31,139.4) node [anchor=north west][inner sep=0.75pt]  [color={rgb, 255:red, 74; green, 144; blue, 226 }  ,opacity=1 ]  {$x,\ y$};
	% Text Node
	\draw (31,173.4) node [anchor=north west][inner sep=0.75pt]  [color={rgb, 255:red, 65; green, 117; blue, 5 }  ,opacity=1 ]  {$x:\ y$};
	% Text Node
	\draw (37,213.4) node [anchor=north west][inner sep=0.75pt]    {$\overline{xy}$};
	% Text Node
	\draw (169,141) node [anchor=north west][inner sep=0.75pt]  [color={rgb, 255:red, 74; green, 144; blue, 226 }  ,opacity=1 ] [align=left] {$\displaystyle x\ \text{Muere}$};
	% Text Node
	\draw (162,177) node [anchor=north west][inner sep=0.75pt]  [color={rgb, 255:red, 65; green, 117; blue, 5 }  ,opacity=1 ] [align=left] {$\displaystyle x:y\ \text{Muere}$};
	% Text Node
	\draw (313,139) node [anchor=north west][inner sep=0.75pt]  [color={rgb, 255:red, 74; green, 144; blue, 226 }  ,opacity=1 ] [align=left] {$\displaystyle y\ \text{Muere}$};
	% Text Node
	\draw (312,207.4) node [anchor=north west][inner sep=0.75pt]    {$\overline{xy} \ \text{Muere}$};
	% Text Node
	\draw (444,170.4) node [anchor=north west][inner sep=0.75pt]  [color={rgb, 255:red, 208; green, 2; blue, 27 }  ,opacity=1 ]  {$T_{x} \ +\ T_{y} \ =\ T_{xy} \ +\ T_{\overline{xy}}$};
	
	
\end{tikzpicture}

$${\color{blue}\Rightarrow \px[t]{x} + \px[t]{y} = \px[t]{xy} + \px[t]{\joint{xy}}}$$

$${\color{red}\Rightarrow \px[t]{\joint{xy}} = \px[t]{x} + \px[t]{y} - \px[t]{xy}}$$


\begin{center}
	(La fórmula anterior es una aplicación del teorema de inclusión exclusión)
\end{center}

\vspace{0.2cm}

\textbf{Ejemplo:}

\vspace{0.2cm}
	
	\[
		\px[2]{5} = 0{.}9 \text{ y } 	\px[3]{77} = 0{.}8
	\]

Para dos variables independientes calcular: 
	
	\begin{enumerate}[label = \alph*)]
		\item $\px[2]{\joint{75:77}} = \px[2]{\joint{75}} + \px[2]{\joint{75}} - \px[2]{75:75} = 2(\px[2]{75}) - (\px[2]{75})^2 = 2(0.9) - (0.9)^2 = 0.99$
		\item $\px[3]{\joint{77:77}} = \px[3]{77} + \px[3]{77} - \px[3]{77:77} = 2(0.8) - (0.8)^2 = 0.96$
		\item $\px[5]{\joint{75:75}} = \px[5]{75} + \px[5]{75} - \px[5]{75:75} = 2(\px[2]{75} \cdot \px[3]{77}) - (\px[2]{75} \cdot \px[3]{77})^2 = 2(0.9 \cdot 0.8) - (0.9 \cdot 0.8)^2 = 0.9216$
	\end{enumerate}

\begin{advertencia}
	
	\textbf{Nota:} es preciso que recuerdes que efectuar la siguiente operación no es válido: 
	
	\[
		\px[5]{\joint{75:75}} = \px[2]{\joint{75:75}} \cdot \px[3]{\joint{77:77}} = 
		0.99 \cdot 0.96 = 
		.9504 
		\neq 0.9216	
	\]
\end{advertencia}

\vspace{0.2cm}

\textbf{Ejercicio:}

\vspace{0.2cm}

\begin{center}
	\begin{tabular}{c|c|c}
		
		$x$ & $\qx{x}^m$ & $\qx{x}^t$  \\
		\hline
		80 & 0.1 & 0.07   \\
		81 & 0.12 & 0.09  \\
		82 & 0.14 & 0.11   \\
		83 & 0.16 & 0.13   \\
		84 & 0.18 & 0.15   \\
		
	\end{tabular}
\end{center}

Juan y laura son pareja, Juan tiene 82, Laura 80 y tienen vidas independientes, se nos pide calcular la probabilidad de que la segunda muerte ocurra durante el 3er año. 



\tikzset{every picture/.style={line width=0.75pt}} %set default line width to 0.75pt        

\vspace{2cm}

\begin{center}
\begin{tikzpicture}[x=0.75pt,y=0.75pt,yscale=-1,xscale=1]
	%uncomment if require: \path (0,300); %set diagram left start at 0, and has height of 300
	
	%Straight Lines [id:da8315110163366394] 
	\draw    (114,141) -- (526.33,140.67) (254,136.89) -- (254,144.89)(394,136.77) -- (394,144.77) ;
	\draw [shift={(526.33,140.67)}, rotate = 179.95] [color={rgb, 255:red, 0; green, 0; blue, 0 }  ][line width=0.75]    (0,5.59) -- (0,-5.59)   ;
	\draw [shift={(114,141)}, rotate = 179.95] [color={rgb, 255:red, 0; green, 0; blue, 0 }  ][line width=0.75]    (0,5.59) -- (0,-5.59)   ;
	%Shape: Rectangle [id:dp8600438542835096] 
	\draw  [color={rgb, 255:red, 255; green, 0; blue, 0 }  ,draw opacity=1 ][fill={rgb, 255:red, 255; green, 0; blue, 0 }  ,fill opacity=0.55 ] (395.33,135.67) -- (528.33,135.67) -- (528.33,146.67) -- (395.33,146.67) -- cycle ;
	%Shape: Brace [id:dp044101165379452745] 
	\draw   (113,184) .. controls (113.01,188.67) and (115.34,191) .. (120.01,190.99) -- (244.18,190.84) .. controls (250.85,190.83) and (254.18,193.16) .. (254.18,197.83) .. controls (254.18,193.16) and (257.51,190.83) .. (264.18,190.82)(261.18,190.82) -- (388.34,190.66) .. controls (393.01,190.66) and (395.34,188.33) .. (395.33,183.66) ;
	
	% Text Node
	\draw (108,152.4) node [anchor=north west][inner sep=0.75pt]    {$0$};
	% Text Node
	\draw (521,151.4) node [anchor=north west][inner sep=0.75pt]    {$3$};
	% Text Node
	\draw (389,150.4) node [anchor=north west][inner sep=0.75pt]    {$2$};
	% Text Node
	\draw (249,150.4) node [anchor=north west][inner sep=0.75pt]    {$1$};
\end{tikzpicture}
\end{center}

\vspace{0.2cm}

$$\qx[2|1]{\joint{82:80}} = \px[2]{\joint{82:80}} - \px[3]{\joint{82:80}} = \underbrace{ \px[2]{82} + \px[2]{80} - \px[2]{82} \cdot \px[2]{80}}_{\px[3]{\joint{82:80}}} - \underbrace{[\px[3]{82} + \px[3]{80} - \px[3]{82} \cdot \px[3]{80} ]}_{\px[2]{\joint{82:80}}}$$

$$\px[2]{82} = (1 - \qx{82}) (1- \qx{83}) = \px{82} \cdot \px{83}$$

$$\px[2]{82} = (1 - 0{.}14) (1 - 0{.}16) = 0{.}7224$$

$$\px[2]{80} = (1 - \qx{80}) (1 - \qx{84})$$

$$(1- 0{.}07)(1 - 0{.}09) = 0{.}8463$$

$$\px[3]{82} = (1 - \qx{82})(1 - \qx{85})(1 - \qx{84}) = (1 - 0{.}14) ( 1 - 0{.}16) ( 1 - 0{.}18) = 0{.}592368$$

$$\px[3]{80} = (1 - \qx{80}) (1 - \qx{81}) (1 - \qx{82}) = (1 - 0{.}07) (1 - 0{.}09) ( 1 - 0{.}11) = 0{.}753207$$

Por lo que: 

$$\qx[2|1]{\joint{82:80}} = 0{.}7224 + 0{.}8463 - (0{.}7224 \cdot 0{.}8463) - [0{.}592368 + 0{.}753207 - (0{.}592368 \cdot 0{.}753207)] = 0{.}0579336042$$

\subsection{Beneficios de vidas conjuntas}

Seguro de vida entera, continuo, vida conjunta

$$\Ax*{xy} = \int_{0}^{\infty} e^{-\delta t} \cdot \px[t]{xy}  \cdot \mu_{xy} dt = \int_{0}^{\infty} e^{-\delta t} \cdot \px[t]{xy} (\mu_x + \mu_y) dt$$

Anualidad conjunta continua

$$\bar{a}_{xy} = \int_0 ^{\infty} e ^{- \delta t} \px[t]{xy} dt$$

\textbf{Recordatorio:} 

$$a_{xy} = \frac{1 - \Ax*{xy}}{\delta}$$

\subsection{Beneficios de último sobreviviente} 

$$\Ax*{\joint{xy}} = \Ax*{x} + \Ax*{y} - \Ax*{xy}$$

$$\bar{a}_{\joint{xy}} = \bar{a}_x + \bar{a}_y - \bar{a}_{xy}$$

$$\bar{a}_{\joint{xy}} = \frac{1 - \Ax*{\joint{xy}}}{\delta}$$

\subsection{Beneficios de temporales} 

$$\Ax*{\tilde{x}:\endow{y}{n}} = \int_0 ^n e^{- \delta t} \px[t]{xy} (\mu_x + \mu_y) dt$$

$$\bar{a}_{\tilde{x}:\endow{y}{n}} = \int_0 ^n e^{- \delta t} \px[t]{xy} dt$$

Seguro ultimo sobreviviente 

$$\Ax*{\tilde{\joint{x:y}}:\actuarialangle{n}} = \bar{a}_{\term{x}{n}} + \bar{a}_{\term{y}{n}} - \bar{a}_{x:\tilde{y}:\actuarialangle{n}}$$

Anualidad ultimo sobreviviente

$$ \bar{A}_{\tilde{\joint{x:y}}:\actuarialangle{n}} $$

Relación: 

$$\bar{a}_{x:y:\actuarialangle{n}} = \frac{1 - \Ax*{xy:\actuarialangle{n}}}{\delta}$$

Relación:
%$$\bar{a}_{\joint{xy}:\actuarialangle{n} = \frac{1 - \Ax*{\joint{xy}:\actuarialangle{n}}}{\delta}$$


\textbf{Ejemplo:}

Sean $(x)$ y $(y)$ dos vidas independientes con fuerzas de mortalidad constantes

\[
	\mu_x = 0{.}002 \quad \mu_y = 0{.}004
\] 

y dada una fuerza de interés

\[
	\delta = 0{.}05
\] 

Se desea estimar el \textbf{Valor Presente Actuarial (VPA)} de un beneficio de \$ 500{,}000 bajo distintos esquemas:

\vspace{0.2cm}

\begin{enumerate}[label =\alph*)]
	\item Vida entera al momento de la primera muerte.
	\item Vida entera al momento de la segunda muerte.
	\item Temporal de 25 años pagadero al momento de la primer muerte. 
	\item Temporal de 25 a;os pagadero al momento de la segunda muerte. 
\end{enumerate}

	\vspace{0.2cm}

	\textbf{Fórmulas bajo fuerza de mortalidad constante}
	
	\vspace{0.2cm}
	
	Para la \textbf{vida conjunta} (primer fallecimiento):
	
	\[
	\mu_{xy} = \mu_x + \mu_y
	\]
	
	\[
	\overline{A}_{xy}
	=
	\frac{\mu_x + \mu_y}{\mu_x + \mu_y + \delta}
	\]
	
	\[
	\overline{A}_{xy:\angl{n}}
	=
	\frac{\mu_x + \mu_y}{\mu_x + \mu_y + \delta}
	\left(
	1 - e^{-n(\mu_x+\mu_y+\delta)}
	\right)
	\]
	
	Para una vida individual:
	
	\[
	\overline{A}_x
	=
	\frac{\mu_x}{\mu_x+\delta}
	\qquad
	\overline{A}_y
	=
	\frac{\mu_y}{\mu_y+\delta}
	\]
	
	Para el \textbf{último sobreviviente} (segunda muerte):
	
	\[
	\overline{A}_{\overline{xy}}
	=
	\overline{A}_x
	+
	\overline{A}_y
	-
	\overline{A}_{xy}
	\]
	
	\vspace{0.5cm}
	
	\textbf{Resolución del ejercicio}
	
	\vspace{0.2cm}
	
	\textbf{a) Vida entera pagadera al momento de la primera muerte}
	
	\[
	\overline{A}_{xy}
	=
	\frac{0.002+0.004}{0.002+0.004+0.05}
	=
	\frac{0.006}{0.056}
	=
	0.1071429
	\]
	
	\[
	VPA
	=
	500{,}000
	\times
	0.1071429
	=
	53{,}571.42
	\]
	
	\textbf{b) Vida entera pagadera al momento de la segunda muerte}
	
	\[
	\overline{A}_{\overline{xy}}
	=
	\overline{A}_x
	+
	\overline{A}_y
	-
	\overline{A}_{xy}
	\]
	
	\[
	=
	\frac{0.002}{0.002+0.05}
	+
	\frac{0.004}{0.004+0.05}
	-
	\frac{0.006}{0.056}
	\]
	
	\[
	=
	0.0384615
	+
	0.0740741
	-
	0.1071429
	=
	0.0053927
	\]
	
	\[
	VPA
	=
	500{,}000
	\times
	0.0053927
	=
	2{,}696.37
	\]
	
	\textbf{c) Seguro temporal 25 años a la primera muerte}
	
	\[
	\overline{A}_{xy:\angl{25}}
	=
	\frac{0.006}{0.056}
	\left(
	1 - e^{-25(0.056)}
	\right)
	\]
	
	\[
	=
	0.1071429
	\left(
	1 - e^{-1.4}
	\right)
	\]
	
	\[
	=
	0.1071429
	(0.7534)
	=
	0.080721
	\]
	
	\[
	VPA
	=
	500{,}000
	\times
	0.080721
	=
	40{,}360.55
	\]
	
	\textbf{d) Seguro temporal 25 años a la segunda muerte}
	
	\[
	\overline{A}_{\overline{xy}:\angl{25}}
	=
	\overline{A}_{x:\angl{25}}
	+
	\overline{A}_{y:\angl{25}}
	-
	\overline{A}_{xy:\angl{25}}
	\]
	
	donde
	
	\[
	\overline{A}_{x:\angl{25}}
	=
	\frac{0.002}{0.052}
	\left(
	1 - e^{-25(0.052)}
	\right)
	\]
	
	\[
	\overline{A}_{y:\angl{25}}
	=
	\frac{0.004}{0.054}
	\left(
	1 - e^{-25(0.054)}
	\right)
	\]
	
	y
	
	\[
	\overline{A}_{xy:\angl{25}}
	=
	\frac{0.006}{0.056}
	\left(
	1 - e^{-25(0.056)}
	\right)
	\]
	
	\vspace{0.5cm}
	
	\textbf{Tarea}
	
	\vspace{0.2cm}
	
	Para un beneficio de \$25{,}000:
	
	\begin{enumerate}[label = \alph*)]
		\item Anualidad vitalicia mientras ambas personas estén con vida.
		\item Anualidad vitalicia mientras al menos una persona esté con vida.
		\item Anualidad temporal 25 años bajo vida conjunta.
		\item Anualidad temporal 25 años bajo último sobreviviente.
	\end{enumerate}
	
	Para el caso discreto
	
	\[
		\Ax{xy} = \sum_{k=0}^\infty v^k \Px[k]{xy}(1 - \Px{x+k:y+k})
	\]

	\[
		\Ax{\joint{xy}} = \sum v^{k+1} \Px[k]{\joint{xy}} (1 - \Px{\joint{x+k:y+k}})
	\]
	
	La relación mencionada anteriormente se respeta 
	
	\[
		\Ax*{\joint{xy}} = \Ax{x} + \Ax{y} - \Ax{xy}
	\]

	Mientras que para los temporales se vería algo así: 
	
	\[
		\Ax{\term{xy}{n}} = \sum_{k=0}^{n-1} v^{k+1} \cdot \qx[k|]{xy}
	\]

	\[
	\Ax{\term{\joint{xy}}{n}} = \sum_{k=0}^{n-1} v^{k+1} \cdot \qx[k|]{\joint{xy}}
	\]
	
	Para las anualidades sería: 
	
	\[
		\ddot{a}_{xy} = \sum_{k=0}^\infty v^k \Px[k]{xy}
	\]
	
	\[
	\ddot{a}_{xy} = \sum_{k=0}^\infty v^k \Px[k]{xy}
	\]
	
	\[
	\ddot{a}_{\endow{x:y}{n}} = \sum_{k=0}^{n-1} v^k \Px[k]{xy}
	\]
	
	\subsection{Valor de rescate}
	
	El monto que la aseguradora le otorga al asegurado en caso de cancelación asumir que la prima es nivelada o única. 
	
	El valor de rescate es la diferencia entre el valor esperado de obligaciones futuras de la aseguradora ej. El valor esperado de obligaciones futuras del asegurado por el pago de primas. 
	
	$$ \Vx[x]{\endow{x}{n}} = \Ax{\endow{x+t}{n-1}} - P \cdot N \hspace{.25cm} _{x:n} \ddot{a}_{\endow{x+t}{n-t}}$$
	
	Forma parte de los valores garantizados cuando el asegurado decide no continuar con el seguro, entonces se devuelve una parte de la reserva, siempre y cuando tenga 2 años. 
	
	En la vida real (Aplicación práctica), las aseguradoras manejan sus propias tablas de cancelación e idealmente cada producto debería tener su propia tabla de cancelación similar a una tabla de mortalidad, con las subsecuentes probabilidades de cancelación. 
	
	\begin{itemize}
		\item Seguro saldado: Si el asegurado decide no seguir pagando su póliza (no más primas) pero desea seguir cubierto. 
		\subitem El valor de rescate de ese momento puede servir para pagar el plazo que falte de transcurrir de la vigencia del contrato. 
		\subsubitem Valor de rescate paga una cobertura a prima única por el mismo plazo pero menor suma asegurada. 
		\subitem Respeta temporalidad. 
		\subitem No más pago de primas. 
		\subitem Menor Suma Asegurada (Lo que alcance con el valor de rescate).
		\item Seguro prorrogado: El valor de rescate sirve para comprar un seguro
		\subitem Misma suma asegurada. 
		\subitem Mejor temporalidad. 
		\subitem Sin pago de primas adicionales. 
	\end{itemize}
	
	\section{Anualidades reversionarias}
	
	Anualidad para una persona y condicionada a la muerte de edad $x$. 
	
	$$ \ddot{a}_{x|y} = \ddot{a}_{y} - \ddot{a}_{xy}$$
	
	$$\ddot{a}_{y|x}  = \ddot{a}_{x} - \ddot{a}_{xy}$$
	
	Lo anterior significa que la anualidad se paga mientras $x$ viva. 
	
	\textbf{Ejercicio}
	
	Una pareja compra un seguro. El esposo tiene 55 años y su esposa 50. Los dos tienen acceso a una pensión. Si el muere antes, ella dejará de recibir dicha pensión. 
	
	El plan consiste en \$20,000 al principio de cada año. Estimar el VPA del beneficio. 
	
	$$M_x = 0.002 M_y = 0.003 \delta = 0.05$$
	
	\section{Decrementos Múltiples}
	
	Es un tipo de modelo en el que se tienen distintas causas por las que un individuo sale de un grupo. 

	\begin{callout-azul}

		\textbf{Ejemplo:}

		Cuando alguna aseguradora necesita distinguir entre los mótivos por los que se está dando de baja a una u a otra persona, ya sea para aplicar exclusiones o beneficios adicionales; ej. Muerte (tanto natural como accidental), invalidez, vejez, etc. 

			\end{callout-azul}
	
	Las causas son denotadas bajo este modelo con variables aleatorias $j$, por lo que ahora en este modelo tendremos ya no solo la variable aleatoría $T(x)$, sino que tendremos dos $T(x)$ y $J$. Por ende bajo este modelo la forma de expresar el tiempo futuro de vida de un modelo de decrementos múltiples en donde $(x)$ está sujeto a todas las posibles causas de salida $\tau$ se expresa como: 

	$$T^{(\tau)}_{(x)}$$
	
	Para esto $j = 1, 2, \dots, m$ son las causas de salida (Total de $m$).
	
	De acuerdo a lo anterior y a la teoría de la probabilidad, podemos afirmar que tenemos entonces una distribución bivariada con por lo menos una de las variables aleatorias discretas $J$
	
	$P(T^{(\tau)}_{(x)} \leq t)$: Probabilidad de que $(x)$ salga o termine antes del tiempo $t$ por cualquier causa. 
	
	$P(J=j) = f_{J}(j)$: Probabilidad de que $(x)$ salga o termine por causa $j$ en cualquier momento. 
	
	% Aquí debería ir la gráfica de la función multivariada
	
	Al integrar, 
	
	$$ 
	P(J=j) = \int_{0}^\infty f_{T^{(\tau)}_{(x)}, J} (t, j)dt
	$$
	
	\[
		= \sum^{m}_{j=1} f_{T^{(\tau)}_{(x)}, J} (t, j)
	\]
	
	Sus respectivas funciones de densidad son: 
	
	Conjunta: $f_{T^{(\tau)}_{(x)}, J} = \px[t]{x}^{(\tau)} \mu_{x+t}^{(\tau)}$
	
	Marginal de t: $f_{T^{(\tau)}_{(x)}}(t) = \px[t]{x}^{(\tau)} \mu_{x+t}^{(\tau)}$
	
	Marginal de j: $P(J=j) = \qx[\infty]{x}^{(j)}$
	
	La condicional de $t$ dado $j$
	
	\[
	f_{T(x)\mid J}(t \mid j)
	=
	\frac{f_{T(x),J}(t,j)}{P(J=j)}
	=
	\frac{{}_t p_x^{(j)} \, \mu_{x+t}^{(j)}}{q_x^{(j)}}
	\]
	
	La condicional de $j$ dado $t$
	
	\[
	f_{J \mid T(x)}(j \mid t)
	=
	\frac{f_{T(x),J}(t,j)}{f_{T(x)}(t)}
	=
	\frac{{}_t p_x^{(j)} \, \mu_{x+t}^{(j)}}{{}_t p_x \, \mu_{x+t}}
	=
	\frac{\mu_{x+t}^{(j)}}{\mu_{x+t}}
	\]
	
	\textbf{Ejemplo}
	
	\[
	f_{T_x^{(r)},J}(t,j) = j e^{-6t}
		=
	\begin{cases}
		e^{-6t}, & j = 1 \\
		2e^{-6t}, & j = 2 \\
		3e^{-6t}, & j = 3
	\end{cases}
	\]
	
	\[
	{}_{4|2} q_x^{(2)} = {}_4 p_x^{(r)} \cdot {}_2 q_{x+4}^{(2)}
	\]
	
	\[
	{}_{4|2} q_x^{(2)} = \int_{4}^{6} f_{T_x^{(r)},J}(t,2)\, dt
	\]
	
	\[
	f_{T_x^{(r)},J}(t,j) = {}_t p_x^{(r)} \, \mu_{x+t}^{(j)} = j e^{-6t}
	\]
	
	Empleando lo anterior, podemos proceder: 
	
	\[
	{}_{4|2} q_x^{(2)} 
	= {}_4 p_x^{(\tau)} \cdot {}_2 q_{x+4}^{(2)}
	= \int_{4}^{6} 2e^{-6t}\, dt
	\]
	
	\[
	= 2 \int_{4}^{6} e^{-6t}\, dt
	= 2 \left[ \frac{e^{-6t}}{-6} \right]_{4}^{6}
	\]
	
	\[
	= -\frac{1}{3} \left( e^{-36} - e^{-24} \right)
	= \frac{e^{-24} - e^{-36}}{3}
	\]
	
	\[
	\approx 0.00002048
	\]
	
	\subsection{Tablas de Mortalidad}
	
	Consideremos un grupo de personas de edad $(x)$ que están sujetas a $(m)$ posibles causas de decremento. 
	
	$$ \lx{x}^{(\tau)} \hspace{1cm} \lx{} \text{\# De personas} \hspace{1cm} \tau \text{: Todas las posibles causas de salida}$$
	
	$$ \dx{x}^{(\tau)} \hspace{1cm} \text{\# de personas que dejan de estar activos por cualquier causa entre edad $x$ y $x+1$}$$
	
	$$ \dx{x}^{(j)} \hspace{1cm} \text{\# de personas que dejan de estar activos por ``j'' causa determinada $j=1, \dots , m$} $$
	
	Recordemos que 
	
	$$ \qx{x} = \frac{dx}{\lx{x}} = 1 - \px{x} = 1 - \frac{\lx{x+1}}{\lx{x}} = \frac{\lx{x} - \lx{x+1}}{\lx{x}}$$
	
	$\px{x} = \frac{\lx{x+1}}{\lx{x}}$ Ahora tenemos $\lx{x}^{(\tau)}, \dx{x}^{(\tau)} y \dx{x}^{(j)}$
	
	$\qx{x}^{(\tau)} = \frac{\dx{x}^{(tau)}}{\lx{x}^{(\tau)}}$ Probabilidad de que salga del grupo en el siguiente año por cualquier causa.
	
	$\qx{x}^{(j)} = \frac{\dx{x}^{(j)}}{\lx{x}^{(\tau)}}$ Probabilidad de que una persona salga del grupo por una específica causa $j$ siendo $j= 1, \dots, m$. 
	
	\textcolor{red}{ Ojo: Las personas estan sujetas a $\tau$ (todas las posibles causas de salida)}
	
	$\dx{x}^{(\tau)} = \sum_{j=1}^m \dx{x}^{(i)}$ Lo que quiere decir que un individuo sólo puede salir por una de todas las posibles causas. 
	
	\textcolor{red}{Total de personas en un año entre $x$ y $x+1$ es la suma de todos los decrementos de las ``$m$'' razones de salida.}
	
	$$\qx[n|]{x}^{(\tau)} = \px[n]{x}^{(\tau)} \cdot \qx{x+n}^{(\tau)} = \frac{\lx{x+n}^{(\tau)}}{\lx{x}^{(\tau)}} \cdot \frac{\dx{x+n}^{(\tau)}}{\lx{x+n}^{(\tau)}} = \frac{\dx{x+n}^{(\tau)}}{\lx{x}^{(\tau)}}$$
	
	Esto es la probabilidad de que $(x)$ salga por cualquier causa entre $x+n$ y $x+n+1$ 
	
	$$\qx[n|]{x}^{(j)} = \px[n]{x}^{(\tau)} \cdot \qx{x+n}^{(j)} = \frac{\lx{x+n}^{(\tau)}}{\lx{x}^{(\tau)}} \cdot \frac{\dx{x+n}^{(j)}}{\lx{x+n}^{(\tau)}} = \frac{\dx{x+n}^{(j)}}{\lx{x}^{(\tau)}}$$
	
	\begin{figure}[h]

	\caption{Tabla a emplear para la siguiente clase}
	\[
	\begin{array}{c|cccccc}
		\text{Año } k & 0 & 1 & 2 & 3 & 4 & 5 \\
		\hline
		q^{(1)}_{x+k} & 0.003 & 0.004 & 0.005 & 0.006 & 0.007 & 0.008 \\
		q^{(2)}_{x+k} & 0.140 & 0.120 & 0.100 & 0.080 & 0.060 & 0.040 \\
		q^{(r)}_{x+k} & 0.143 & 0.124 & 0.105 & 0.086 & 0.067 & 0.048 \\
	\end{array}
	\]
	\end{figure}
	
\end{document}          
}